\uuid{MtAc}
\titre{Vecteurs}
\auteur{bodin ; hanani ; mzari}
\organisation{exo7}
\licence{CC BY-NC-SA 4.0}
\datecreate{2020-01-20}
\niveau{L1}
\module{Géométrie}
\chapitre{Géométrie du plan}
\sousChapitre{Distance, norme, produit scalaire}
\difficulte{1}
\theme{vecteurs}
\exo7qcmid{lille-1.geometrie-plane.003}

\contenu{
\begin{question}
\qclassification{140.01,140.02}
\qtags{motcle=vecteurs}
On considère les points $A(1,1)$, $B(-1,1)$ et $C(1,-1)$. Quelles sont les bonnes réponses ?
\begin{answers}
    \bad{Les vecteurs $\overrightarrow{AB}$ et $\overrightarrow{AC}$ sont égaux.}
    \bad{$\overrightarrow{AB}=-\overrightarrow{AC}$}
    \bad{Les vecteurs $\overrightarrow{AB}$ et $\overrightarrow{AC}$ sont colinéaires.}
    \good{Les vecteurs $\overrightarrow{AB}$ et $\overrightarrow{AC}$ sont orthogonaux.}
\end{answers}
\begin{explanations}
D'abord, $\overrightarrow{AB}=(-2,0)$ et $\overrightarrow{AC}=(0,-2)$, et puis le produit scalaire $\overrightarrow{AB}\cdot\overrightarrow{AC}$ est nul. Donc $\overrightarrow{AB}$ et $\overrightarrow{AC}$ sont orthogonaux.
\end{explanations}
\end{question}
}
